\scanpage{169}
% L1–7
Considérons maintenant
\[
\underset{\text{sous }L_1}{(A_1,\sigma_1)}\quad
\underset{\text{sous }L_2}{(A_2,\sigma_{\unsure{2}})}
\]
droites affines à involution, d'où
\[
\underset{\text{\unsure{sous|sur} }L_1\otimes L_2}{A=\deleted{Q(A_1}A_1*A_2}
\quad\deleted{\unsure{b_1}}
\]
je dis que
\[
\underbrace{Q(A)}_{\text{sous }L^{\otimes 2}}\simeq
\underbrace{Q(A_1)}_{\text{sous }L_1^{\otimes 2}}
\deleted{\otimes}{*}\underbrace{\deleted{\unsure{A}}Q(A_2)}_{\text{sous }L_2^{\otimes 2}}
\]
\margin{??}
(iso canonique)

% L8–10
\margin{?}Pour donner un sens à cette
formule, il faut définir sur un
$Q(A)$ une involution naturelle

% L11–17
En tous cas, \deleted{en termes} on a
une application canonique
\[
\begin{aligned}
&q_1,q_2\mathbin{\text{---}}q_1\otimes q_2\mid \unsure{\check A_1|\dot A_1|A_1}*A_2\subset E_1\otimes E_2\\
&Q(A_1)\times Q(A_2)\longrightarrow Q(A)
\end{aligned}
\]
\deleted{\unsure{donnée en termes de co}}
qui s'explicite ainsi, si on choisit
\deleted{\unsure{pour|base}} pt base $o_1\in A_1$, $o_2\in A_2$, de sorte
que
\[
A_i\simeq A_i^0,\quad \sigma_i\text{ devient }x\mathrel{\vdash\!\!-\!\!-}b_i-x
\]
% L18–21
les formes quadratiques sur $A_i$ \deleted{\unsure{dev|comp}}compatibles
à $\sigma_i$ sont
\[
q_{c_i}\text{ ou }q_{b_i,c_i}\quad x_i\longmapsto x_i^2-\deleted{\unsure{t_i}}b_i\unsure{x.|x_i}+c_i
\qquad c_i\in L^{\otimes 2}\quad\unsure{avec|conv.|comme}
\]
et on a alors
\[
q_{c_{\unsure{i|1}}}\otimes q_{\unsure{c_2|b_2}}=q_c\quad
\underset{\substack{\rotatebox{-90}{$\in$}\\L_1\otimes L_2}}{x}
\mathrel{\vdash\!\!-\!\!-}x_{\deleted{\unsure{12}}}^{2}-(b_1b_2)x_{\unsure{2|{}}}+
(b_2^2c_1+b_1^2c_2-4c_1c_2)
\]
% archive sheet number: 168
